% ECERC Revisited: Final Report
% Compile with pdfLaTeX (twice, so the long tables settle). Keep the figures/ folder next to this file.
% The figures are vector PDFs converted from the SVG files in ../architecture/.
\documentclass[11pt,a4paper]{article}

\usepackage[T1]{fontenc}
\usepackage[utf8]{inputenc}
\usepackage{lmodern}
\usepackage{textcomp}
\usepackage[margin=2cm]{geometry}
\usepackage{amsmath,amssymb}
\usepackage{graphicx}
\usepackage{float}
\usepackage{array}
\usepackage{booktabs}
\usepackage{longtable}
\usepackage{adjustbox}
\usepackage{enumitem}
\usepackage[hyphens]{url}
\usepackage[hidelinks]{hyperref}

\setcounter{secnumdepth}{0}
\setlength{\parskip}{0.55em}
\setlength{\parindent}{0pt}
\setlength{\tabcolsep}{4pt}
\setlength{\LTpre}{0.6em}
\setlength{\LTpost}{0.6em}
\renewcommand{\arraystretch}{1.15}
\setlist[itemize]{topsep=0.2em,itemsep=0.2em}

\newcolumntype{L}[1]{>{\raggedright\arraybackslash}p{#1}}
\newcommand{\mn}{\textminus}
\newcommand{\archfig}[2]{%
  \begin{figure}[H]
    \centering
    \includegraphics[width=\linewidth,height=0.78\textheight,keepaspectratio]{figures/#1.pdf}
    \caption{#2}
  \end{figure}}

\title{ECERC Revisited: Final Report}
\author{Divyansh Dubey\\12340690 \and Keshav Mishra\\12341140 \and Lakshay Gupta\\12341300}
\date{Oct 30, 2026}

\begin{document}

\maketitle

\section{Summary}

We extended ECERC on three limitations and beat the authors' checkpoints on both datasets, in two separate ways. The final systems reach 72.43 to 72.77 weighted F1 (wF1) on IEMOCAP against 71.78, and 53.77 to 53.82 macro-F1 (mF1) on MELD against 50.87.

The team built two merged systems, and they agree on what works:

\begin{itemize}
  \item \textbf{Approach A, merged per dataset:} each teammate merged Confidence-Aware Cause Gating (CACG) into their own pipeline and retrained on their own dataset.
  \item \textbf{Approach B, one codebase:} all three limitations were merged into one rewritten codebase and retrained on a laptop.
\end{itemize}

\begin{center}
\begin{adjustbox}{max width=\linewidth}
\begin{tabular}{lccc}
\toprule
Dataset & Authors' checkpoints & Approach A & Approach B \\
\midrule
IEMOCAP, ensemble wF1 & 71.78 & 72.43 & 72.77 \\
IEMOCAP, ensemble mF1 & 70.82 & 71.14 & 71.52 \\
MELD, wF1 & 66.46 & 67.11 & 67.10 \\
MELD, mF1 & 50.87 & 53.77 & 53.82 \\
\bottomrule
\end{tabular}
\end{adjustbox}
\end{center}

Better inputs produced the accuracy gains, not new losses. WavLM audio drove the IEMOCAP gain, and zero-shot LLM fusion drove the MELD gain. Across Limitations 1 and 2, 17 earlier methods fell short of their goal, and those failures pointed to the methods that worked. CACG added no proven accuracy, but it makes cause weighting readable.

The gaps between Approaches A and B are too small to rank them. Approach A's IEMOCAP result has stronger evidence, with 20 seeds against 5.

\section{Project overview}

The team kept ECERC's structure and addressed one failure mode per person. ECERC (Zhang and Tan, ACL 2025) recognises emotions in conversation by combining text, audio and visual evidence with four causes: the speaker's own past emotion, other speakers' emotions, the utterance's own event, and other speakers' events.

\begin{center}
\begin{tabular}{L{3.6cm}L{2.7cm}L{1.6cm}L{7.8cm}}
\toprule
Limitation & Owner & Dataset & Problem \\
\midrule
L1: category confusion & Divyansh Dubey & IEMOCAP & Happy is confused with Excited, and Angry with Frustrated \\
L2: class imbalance & Keshav Mishra & MELD & Fear and Disgust are under 6\% of the data and score poorly \\
L3: cause gating & Lakshay Gupta & Both & Causes are weighted the same way whatever the evidence, and predictions are over-confident \\
\bottomrule
\end{tabular}
\end{center}

\archfig{overview}{How the individual solutions feed the two merged systems}

\textbf{Two approaches.} In Approach A, each dataset was merged separately. Divyansh added CACG to the authors' IEMOCAP code, and Keshav added L2 to a copy of the rewritten code from L3. In Approach B, all three limitations sit in one rewritten codebase behind command-line flags, and every run was retrained on one laptop.

\textbf{Datasets.} IEMOCAP uses Sessions 1 to 4 for training and Session 5 for testing, with 1,623 test utterances and 6 emotions. The first 10\% of training conversations form the validation set. MELD has 7 emotions:

\begin{center}
\begin{tabular}{lcccc}
\toprule
MELD split & Dialogues & Utterances & Fear & Disgust \\
\midrule
Train & 1,038 & 9,989 & 268 & 271 \\
Dev & 114 & 1,109 & 40 & 22 \\
Test & 280 & 2,610 & 50 & 68 \\
\bottomrule
\end{tabular}
\end{center}

\textbf{Metrics.} Weighted F1 (wF1) is the paper's main metric, and macro-F1 (mF1) gives every emotion equal weight. Expected calibration error (ECE) is the gap between confidence and accuracy, in percent, and lower is better. TS means temperature scaling fitted on validation data. Pair errors count confusions between two similar emotions.

\section{Baseline and reproduction}

We reproduced the authors' results with their checkpoints, but no retraining reached them. The rewrite was verified as equivalent to the authors' model, so it is a correct base to build on, not an accuracy improvement.

\archfig{ecerc_baseline}{ECERC as released by the authors}

\begin{center}
\begin{adjustbox}{max width=\linewidth}
\begin{tabular}{lllcccc}
\toprule
Source & Code & Runs & IEMOCAP acc & IEMOCAP wF1 & MELD acc & MELD wF1 \\
\midrule
ECERC paper & Authors' & Mean of 3 & 71.60 & 71.78 & 67.32 & 66.46 \\
Authors' checkpoints & Authors' & Mean of 3 & 71.60 & 71.78 & 67.32 & 66.46 \\
One authors' checkpoint & Rewritten & 1 & 71.60 & 71.77 & 67.55 & 66.58 \\
Retrained, Divyansh Dubey & Authors' & 20 seeds & 69.50 & 69.70 & & \\
Retrained, Keshav Mishra & Authors' & 3 seeds & & & 66.87 & 66.00 \\
Retrained, Keshav Mishra & Rewritten copy & 3 seeds & & & 66.83 & 65.72 \\
Retrained, Lakshay Gupta & Rewritten & 5 and 3 seeds & 69.34 & 69.52 & 67.08 & 65.90 \\
\bottomrule
\end{tabular}
\end{adjustbox}
\end{center}

A blank cell means that person did not run that dataset.

\textbf{The rewrite.} The rewritten code puts both datasets in one package, keeps the authors' parameter names and adds a flag for every extension. With the same weights and inputs, it gives the same outputs as the official model, with a maximum difference of 0.0. The authors' checkpoints load into it unchanged. It changes how the code is organised, not what the model computes.

\textbf{Why retraining scores lower.} IEMOCAP retraining lands about 2 wF1 below the checkpoints in both code bases, so the rewrite is not the cause. The authors say exact reproduction needs their environment. On IEMOCAP, seeds alone move wF1 by over 2 points (68.45 to 70.79 across 20 seeds). The released IEMOCAP checkpoints all score 71.77 to 71.78, so they are likely the best of several runs.

Because of this gap, each limitation compares against a retrained baseline, and the checkpoints serve as a reference for overall accuracy.

\textbf{Problems found in the released code.} We kept all three, so that the baseline stays comparable with the authors' model.

\begin{itemize}
  \item \textbf{No cause-level softmax.} The project proposal described a softmax over the four causes, but the release uses an element-wise sigmoid gate on each one.
  \item \textbf{Empty-history attention.} A cause with no valid history gets uniform attention over the whole padded dialogue, including future turns. This affects 28.6\% (same speaker) and 15.8\% (other speakers) of MELD test utterances, and 3 to 4\% on IEMOCAP.
  \item \textbf{MELD modality routing.} The MELD loader swaps the audio and visual slots.
\end{itemize}

\section{Limitation 1: category confusion (Divyansh Dubey, IEMOCAP)}

Five training changes on the original features failed to fix the confused pairs, which showed the real problem was missing information. Better audio then worked: WavLM with VAD supervision raised wF1 by 1.43 and removed 20\% of clear-cut Angry--Frustrated errors. All L1 runs used the authors' code and were compared with the same 20 seeds, paired seed by seed.

\subsection{Every method tried}

{\footnotesize
\begin{longtable}{L{0.4cm}L{2.4cm}L{2.6cm}L{3.9cm}L{1.8cm}L{3.9cm}}
\toprule
\# & Method & Why it was tried & Result & Verdict & Why \\
\midrule
\endhead
\bottomrule
\endlastfoot
1 & Confusion-aware contrastive loss & Separate the internal vectors of confused emotions & Looked like about \mn15 pair errors on 5 seeds; 174.3 against 173.9 on 10 seeds, wF1 69.63 against 69.79 & Failed & The classes are not separated even on training data, and the 5-seed gain was noise \\
2 & Pair penalty & Punish probability on the partner class & Same 10-seed result, alone or combined & Failed & It can only move the boundary. Even a shift tuned on test added just 1.2 wF1, and tuned on validation it lost 1.5 \\
3 & Ensembling & Cancel random errors & wF1 69.70 to 70.26 and ECE 12.0 to 10.0 with 5 models, but pair errors stayed & Helped overall, not the pairs & Different runs get the same utterances wrong \\
4 & VAD supervision, original audio & The pairs differ in energy and dominance & Angry--Frustrated errors 97.9 to 95.9 (15 of 20 runs); wF1 unchanged at 69.78 & Small real gain & It learned arousal (correlation 0.73), but the old audio holds little about how something is said \\
5 & Soft labels from annotator votes & Stop punishing agreement with minority annotators & ECE 12.1 to 6.1, NLL 0.84 to 0.74, wF1 69.82, Happy--Excited errors up about 5 & Calibration only; dropped & The model learned to hedge, and hedges tipped towards the larger class. Temperature scaling calibrates without this cost \\
6 & WavLM audio + VAD supervision & Add new information to the model & wF1 +1.43, Angry--Frustrated clear errors down 20\% & Final method & Modern speech features carry what the old ones lacked \\
\end{longtable}
}

\subsection{The final method}

WavLM-large (layer 20, averaged over time) gives a 1,024-d vector that replaces the 1,582-d IS10 audio features for all 7,433 utterances. Text, video and the model stay the same. A training-only head predicts the human valence, arousal and dominance ratings with loss weight 1.0.

\[
L = w_y \, \mathrm{CE}(p, y) + \lambda_{vad}\, \mathrm{MSE}\big(\mathrm{head}(R),\ \mathrm{zscore}(\mathrm{VAD})\big), \qquad \lambda_{vad} = 1
\]

An audio-only check confirmed the potential before full training, and it predicted the final pattern.

\begin{center}
\begin{tabular}{lcc}
\toprule
Audio-only classifier & IS10 & WavLM \\
\midrule
6-class macro-F1 & 34.8 & 51.8 \\
Angry vs Frustrated, balanced accuracy & 60.5\% & 72.8\% \\
Happy vs Excited, balanced accuracy & 64.2\% & 66.5\% \\
\bottomrule
\end{tabular}
\end{center}

\archfig{limitation1_wavlm_vad}{Limitation 1: WavLM audio and a training-only VAD head}

\subsection{Results}

\begin{center}
\begin{adjustbox}{max width=\linewidth}
\begin{tabular}{lcccc}
\toprule
Metric (20 paired runs) & Baseline & WavLM + VAD & Change [95\% CI] & Runs better \\
\midrule
Accuracy & 69.50 & 70.99 & +1.48 [+1.06, +1.97] & 19 of 20 \\
wF1 & 69.70 & 71.13 & +1.43 [+1.00, +1.87] & 19 of 20 \\
mF1 & 68.83 & 69.73 & +0.90 [+0.47, +1.35] & 15 of 20 \\
ECE after TS & 4.30 & 3.40 & \mn0.90 [\mn1.33, \mn0.47] & 16 of 20 \\
Angry--Frustrated clear errors & 16.9 & 13.5 & \mn3.4 [\mn5.1, \mn1.7] & 16 of 20 \\
Happy--Excited clear errors & 12.4 & 12.5 & +0.1 [\mn1.2, +1.5] & 7 of 20, 3 tied \\
\bottomrule
\end{tabular}
\end{adjustbox}
\end{center}

Clear errors count utterances where every annotator voted for the gold label. On the rest, the annotators disagreed, so the correct label is debatable. WavLM alone gives +1.28 wF1, and VAD adds a further +0.15, which is within noise.

Happy--Excited stays unsolved. The two sound almost the same, and annotators drew the line between them inconsistently.

\section{Limitation 2: class imbalance (Keshav Mishra, MELD)}

Twelve methods gave no confirmed macro-F1 gain, which showed the limit is ECERC's evidence, not its loss, decision rule or training budget. Zero-shot LLM fusion then raised test macro-F1 by 3.38 on the authors' checkpoints (95\% CI +0.64 to +5.73). Each reason below is an interpretation of the results, not a proven cause.

\subsection{Every method tried}

{\footnotesize
\begin{longtable}{L{2cm}L{3.5cm}L{3.4cm}L{2.1cm}L{4.5cm}}
\toprule
Group & Method & Result & Verdict & Likely reason \\
\midrule
\endhead
\bottomrule
\endlastfoot
Loss and sampling & Class-balanced loss (CB0.999) & Test mF1 49.52 against 50.40 & Failed & Fear false positives rose from 25 to 95 per run; the extra Fear predictions were only 9\% correct \\
Loss and sampling & Stronger focal loss ($\gamma$ = 2) & 50.06 & Failed & The authors' loss is already focal, so this is a small change \\
Loss and sampling & Oversampling Fear and Disgust dialogues & 50.02 & Failed & It repeats the same few examples and adds no new evidence \\
Loss and sampling & Focal $\gamma$ = 2 + oversampling & 49.13, the lowest & Failed & It combines both weaknesses \\
Decision rule & Post-hoc logit adjustment & \mn0.24 [\mn1.06, +0.68] & Failed & It moved mostly correct predictions: 19 extra Fear predictions gave 1 correct one \\
Representation & Mean of 4 RoBERTa layers & Dev \mn0.07 & Failed the gate & Fear and Disgust stayed at 30 to 36 dev F1 under every candidate \\
Representation & Confusion-aware contrastive loss & Dev \mn0.22 & Failed the gate & Overlapping emotions stay hard to separate, whatever the loss \\
Representation & Auxiliary sentiment head & Dev +0.19 & Failed the gate & Below the +0.5 needed to reach test \\
Representation & Learning rate 3e-5 & Dev \mn0.06 & Failed the gate & It peaked 20 epochs earlier then overfit, so under-training is not the cause \\
Trained expert & RoBERTa-large text expert & +0.27 [\mn1.92, +2.29] & Inconclusive & Fear rose 3.40, but Disgust and Joy fell \\
Trained expert & Qwen2.5-7B LLM, fine-tuned & Cross-fitted dev \mn0.98 & Failed the gate & It learned MELD's skew, and reached only 22 to 23 dev Disgust F1 \\
Better evidence & Whisper-large-v3 audio & Dev +0.08 & Failed the gate & Text dominates on MELD; better audio moved Sadness up and Joy down \\
Final & \textbf{Zero-shot LLM fusion} & \textbf{+3.38 [+0.64, +5.73]} & \textbf{Worked, with a trade-off} & It was never trained on MELD's skewed labels, so its errors differ from ECERC's \\
After final & Distilling the LLM into ECERC & Cross-fitted dev \mn0.42 & Failed the gate & The LLM is still needed at inference \\
\end{longtable}
}

The first group's test baseline is 50.40 mF1. Every variant there was below it on all three seeds. A ``gate'' is a preset minimum gain on dev data that a method needed before it could touch the test set.

\subsection{The final method}

Qwen2.5-32B-Instruct (4-bit, frozen) scores every utterance from the dialogue so far. It is never trained on MELD. Its scores are calibrated without labels, using its mean dev prediction $q$ and the training prior $\pi$, and then combined with ECERC's after training. The weight $w = 0.2$ was chosen on dev.

\[
\tilde{p}_{LLM}(c) \propto \frac{p_{LLM}(c)}{q(c)} \, \pi_{train}(c)
\qquad
\log p_{final} = (1-w) \log p_{ECERC} + w \log \tilde{p}_{LLM}
\]

\archfig{limitation2_llm_fusion}{Limitation 2: ECERC with zero-shot LLM fusion}

The cross-fitted dev gain was +5.63 against a +0.5 gate, so fusion went to the test set once. It works because the LLM matches ECERC on macro-F1 with opposite errors: stronger on Disgust (37.58 against 31.43) and weaker on Joy (55.57 against 64.95).

\subsection{Results}

\begin{center}
\begin{adjustbox}{max width=\linewidth}
\begin{tabular}{lcccccc}
\toprule
System & Acc & wF1 & mF1 & Fear F1 & Disgust F1 & Joy F1 \\
\midrule
Authors' checkpoints & 67.32 & 66.46 & 50.87 & 26.12 & 31.43 & 64.95 \\
Authors' checkpoints + LLM & 68.24 & 67.26 & 54.25 & 33.77 & 40.22 & 61.67 \\
LLM alone & 63.10 & 63.14 & 51.05 & 31.88 & 37.58 & 55.57 \\
Retrained baseline & 66.83 & 65.72 & 49.57 & 23.21 & 30.46 & 64.18 \\
Retrained baseline + LLM & 68.11 & 67.01 & 53.59 & 32.31 & 38.85 & 61.62 \\
\bottomrule
\end{tabular}
\end{adjustbox}
\end{center}

Fusion gained 4.02 mF1 on the retrained baseline (95\% CI +1.55 to +6.45). The pre-registered verdict is a trade-off, not a clean success, because Joy F1 fell 3.28 points, past the \mn2.0 limit set before the run.

\section{Limitation 3: cause gating and calibration (Lakshay Gupta, both datasets)}

CACG did not change accuracy measurably, but it makes cause weighting readable and it slightly lowers calibration error. Temperature scaling gave the largest calibration gain, cutting IEMOCAP ECE from 14.36 to 3.72, and it works on any model.

\subsection{Method}

CACG sits after ECERC's Feature Gating and adds a distribution over the four causes for each utterance. Its parameters add 0.11\% to ECERC's size. With the mask off, it starts identical to ECERC, which the tests confirm.

\archfig{limitation3_cacg}{Limitation 3: the CACG cause gate and temperature scaling}

\textbf{Evidence descriptor.} A 12-d summary holds 3 modality-gate means, 3 history counts and 6 pairwise cosines between cause vectors.

\begin{multline*}
h_i = \bigl[\, \bar{w}^{t}_i,\ \bar{w}^{a}_i,\ \bar{w}^{v}_i,\ \log(1+n_i),\ \log(1+n^{same}_i),\\
\log(1+n^{other}_i),\ \cos(\tilde{f}_{ia}, \tilde{f}_{ib})_{a<b} \,\bigr]
\end{multline*}

\textbf{Temperature and masked weights.} A small network predicts $\tau$ from $h$. A linear scorer gives one logit per cause, and $m$ removes causes with no history.

\[
\tau_i = 0.25 + 3.75 \cdot \sigma\big(\mathrm{MLP}(h_i)\big)
\qquad
g_{ik} = \frac{m_{ik} \exp(z_{ik} / \tau_i)}{\sum_{j} m_{ij} \exp(z_{ij} / \tau_i)}
\]

\textbf{Fusion.} Multiplying by 4 means a uniform gate leaves the features unchanged.

\[
R_i = [\, 4 g_{i1} \tilde{f}_{i1};\ 4 g_{i2} \tilde{f}_{i2};\ 4 g_{i3} \tilde{f}_{i3};\ 4 g_{i4} \tilde{f}_{i4} \,]
\]

\textbf{Utility gate loss (optional).} Each cause gets an auxiliary classifier. The target $q$ says which cause best explains the true emotion, so the gate is trained against the label.

\[
q_{ik} = \frac{m_{ik}\, p_k(y_i \mid \tilde{f}_{ik})^{1/T}}{\sum_{j} m_{ij}\, p_j(y_i \mid \tilde{f}_{ij})^{1/T}}
\qquad
L_{total} = L_{task} + \lambda_{gate}\, \mathrm{KL}(q_i \,\|\, g_i) + \lambda_{aux} L_{aux}
\]

\textbf{Temperature scaling.} One number $T$ divides the logits. It is fitted on validation data only, so it never changes which emotion is predicted.

\[
\hat{p}_i = \mathrm{softmax}(l_i / T),
\qquad
T^{*} = \arg\min_T \sum_{i \in \mathrm{valid}} -\log \hat{p}_{i, y_i}
\]

\subsection{Every method tried}

{\footnotesize
\begin{longtable}{L{0.4cm}L{2.5cm}L{2.6cm}L{3.8cm}L{1.4cm}L{4.3cm}}
\toprule
\# & Method & Why it was tried & Result & Verdict & Why, and what came next \\
\midrule
\endhead
\bottomrule
\endlastfoot
1 & Empty-history attention fix & Stop attention spreading over future turns & IEMOCAP wF1 70.11 to 68.99; MELD 65.39 to 65.36 & Failed & It adds a fake history signal; replaced by the mask \\
2 & Masked cause softmax, fixed $\tau$ & Create the missing cause distribution & IEMOCAP 69.78, MELD 65.32 & Neutral & It gives cause weights but no accuracy; moved to per-utterance $\tau$ \\
3 & Dynamic $\tau$, no gate loss & Let weak evidence soften the weights & IEMOCAP 69.77 on 5 seeds; $\tau$ spread only 0.008 & Kept & Nothing in training rewards $\tau$ for varying, so it stays nearly fixed \\
4 & Proposal's entropy loss & Train the gate as the proposal planned & IEMOCAP 69.36; MELD 65.20; Disgust F1 30.2 to 21.0 & Failed & Its target comes from the same inputs as $\tau$ and never checks the label; replaced by the utility loss \\
5 & Utility loss, $T$ = 1 & Ground the gate target in the label & MELD 65.51 against 65.90 at $\gamma$ = 1; flattens the weights & No gain & The small gain seen at $\gamma$ = 0 did not hold \\
6 & Utility loss, $T$ = 0.25 & Test a sharper target & IEMOCAP 69.52, MELD 65.50 & No gain & Dropped \\
7 & Global temperature scaling & Fix over-confidence directly & IEMOCAP ECE 14.36 to 3.72, predictions unchanged & Worked & Kept in every final system \\
8 & Cause-aware temperature scaling & Set $T$ per utterance from gate signals & 4.11 against 2.95 on IEMOCAP & Failed & Over-confidence is about the same at every context length, so one $T$ already fits \\
9 & Availability mask & Remove causes with no history & Applies to 5.2\% of IEMOCAP and 33.7\% of MELD test utterances & Kept & Removes noisy vectors; effect not measured alone \\
10 & Seed ensemble & Reduce seed noise without retraining & IEMOCAP CACG 69.77 to 70.68 wF1 & Worked & Kept \\
11 & Selective prediction & Hold back risky answers & Accuracy about 70\% to 84.5\% on the most confident half & Worked & Became the demo's low-confidence flag \\
\end{longtable}
}

Rows 1, 2, 4, 6 and 8 used 3 seeds, and their MELD figures were trained at $\gamma$ = 0 before we found the release uses $\gamma$ = 1. These explanations are our interpretation, and we did not test them separately.

\subsection{Results}

\begin{center}
\begin{adjustbox}{max width=\linewidth}
\begin{tabular}{llcllcc}
\toprule
Dataset and run & Model & Seeds & wF1 & mF1 & ECE raw & ECE + TS \\
\midrule
IEMOCAP, authors' code & ECERC & 20 & 69.70 & 68.83 & & 4.30 \\
IEMOCAP, authors' code & CACG & 20 & 69.58 & & & 3.74 \\
IEMOCAP, rewritten code & ECERC & 5 & 69.52 $\pm$ 0.97 & 68.75 & 13.33 & 4.55 \\
IEMOCAP, rewritten code & CACG, no gate loss & 5 & 69.77 $\pm$ 1.28 & 69.06 & 14.36 & 3.72 \\
MELD, rewritten code & ECERC & 3 & 65.90 $\pm$ 0.28 & 49.78 & 3.62 & 2.92 \\
MELD, rewritten code & CACG, utility loss & 3 & 65.51 $\pm$ 0.08 & 49.25 & 3.97 & 3.20 \\
MELD, rewritten copy (Keshav Mishra) & ECERC & 3 & 65.72 $\pm$ 0.38 & 49.57 & 3.02 & 3.49 \\
MELD, rewritten copy (Keshav Mishra) & CACG, no gate loss & 3 & 65.96 $\pm$ 0.24 & 49.67 & 3.90 & 3.05 \\
\bottomrule
\end{tabular}
\end{adjustbox}
\end{center}

A blank cell means the source does not report that value. CACG changes wF1 by \mn0.39 to +0.25 across these runs, which is inside seed noise. The strongest calibration evidence is Divyansh's 20-seed rerun, where CACG lowered scaled ECE in 15 of 20 runs.

\begin{center}
\begin{adjustbox}{max width=\linewidth}
\begin{tabular}{lcccc}
\toprule
Ablation grid, 3 seeds & IEMOCAP wF1 & IEMOCAP ECE + TS & MELD wF1 ($\gamma$ = 0) & MELD ECE + TS ($\gamma$ = 0) \\
\midrule
ECERC & 70.11 $\pm$ 0.46 & 4.21 & 65.39 $\pm$ 0.34 & 3.01 \\
ECERC + attention fix & 68.99 $\pm$ 1.36 & 5.66 & 65.36 $\pm$ 0.22 & 3.18 \\
Masked softmax, fixed $\tau$ & 69.78 $\pm$ 1.11 & 3.87 & 65.32 $\pm$ 0.15 & 2.41 \\
Dynamic $\tau$, no gate loss & 70.65 $\pm$ 0.86 & 2.95 & 65.34 $\pm$ 0.16 & 2.42 \\
Dynamic $\tau$ + entropy loss & 69.36 $\pm$ 1.29 & 4.26 & 65.20 $\pm$ 0.18 & 2.54 \\
Dynamic $\tau$ + utility loss, $T$ = 1 & 69.44 $\pm$ 1.03 & 4.22 & 65.59 $\pm$ 0.04 & 3.18 \\
Dynamic $\tau$ + utility loss, $T$ = 0.25 & 69.52 $\pm$ 0.32 & 3.55 & 65.50 $\pm$ 0.13 & 2.91 \\
\bottomrule
\end{tabular}
\end{adjustbox}
\end{center}

The +0.54 wF1 for dynamic $\tau$ on 3 IEMOCAP seeds shrank to +0.25 on 5 seeds and \mn0.12 on 20, so these 3-seed differences should not be read as gains.

\subsection{Interpretability}

CACG shows which of the four causes the model relied on for each utterance, which ECERC cannot do directly. ECERC's gates act inside each cause vector, so they give no single score comparing one cause with another.

\begin{center}
\begin{adjustbox}{max width=\linewidth}
\begin{tabular}{lcccc}
\toprule
System & Self-contagion & Cross-emotion & Self-event & Cross-event \\
\midrule
IEMOCAP, WavLM + CACG (Approach B) & 0.50 & 0.38 & 0.08 & 0.05 \\
IEMOCAP, WavLM + VAD + CACG (Approach A) & 0.48 & 0.38 & 0.09 & 0.06 \\
MELD, no gate loss, full history (Approach A) & 0.34 & 0.42 & 0.14 & 0.10 \\
MELD, utility loss, full history (Approach B) & 0.28 & 0.27 & 0.23 & 0.23 \\
\bottomrule
\end{tabular}
\end{adjustbox}
\end{center}

\begin{itemize}
  \item \textbf{Emotion causes dominate.} On IEMOCAP they carry 88\% of the weight in the Approach B model.
  \item \textbf{Emotions differ.} Anger leans towards the speaker's own past (0.62 against 0.24), while Sad and Neutral draw about evenly on both speakers.
  \item \textbf{History matters.} Self-contagion grows from 0.39 with 1 to 2 prior utterances to 0.48 after more than 20. With no history, all weight goes to self-event.
  \item \textbf{The utility loss flattens the weights.} On MELD with full history, its gate entropy is 1.38 against a maximum of 1.39, so its weights say little. Without a gate loss, the weights stay readable.
\end{itemize}

These weights have clear limits. They show what the model relied on, not what truly caused an emotion, because we had no human cause labels to check them against. With full history, all seeds agree on the top cause for only 59.5\% of IEMOCAP utterances. Gate entropy is almost the same on right and wrong predictions, so the weights do not signal errors. $\tau$ also stayed nearly constant in every variant without a gate loss.

\section{Approach A: merged per dataset}

Each teammate merged CACG into their own pipeline and reported one dataset. Both systems beat the authors' checkpoints on overall score, and neither gain comes from CACG itself.

\subsection{IEMOCAP: L1 + L3 (Divyansh Dubey)}

This system uses WavLM audio, VAD supervision, CACG with no gate loss, and temperature scaling, built on the authors' code. It was retrained from scratch on 20 seeds.

\archfig{approach_a_iemocap}{Approach A on IEMOCAP: WavLM, VAD head, CACG, temperature scaling and a 5-run ensemble}

\begin{center}
\begin{adjustbox}{max width=\linewidth}
\begin{tabular}{lccccc}
\toprule
Model & Acc & wF1 & mF1 & ECE raw & ECE + TS \\
\midrule
Authors' checkpoint 1 & 71.60 & 71.77 & 70.73 & 11.12 & 2.52 \\
Authors' checkpoint 2 & 71.53 & 71.78 & 70.70 & 11.07 & 2.95 \\
Authors' checkpoint 3 & 71.66 & 71.78 & 71.02 & 11.30 & 2.97 \\
Authors' checkpoints, average & 71.60 & 71.78 & 70.82 & 11.16 & 2.81 \\
Merged, single run (mean of 20) & 71.36 & 71.59 & 70.19 & 13.64 & 3.41 \\
\textbf{Merged, 5-run ensemble} & \textbf{72.18} & \textbf{72.43} & \textbf{71.14} & 10.60 & \textbf{2.12} \\
\bottomrule
\end{tabular}
\end{adjustbox}
\end{center}

The ensemble row is the mean of 4 separate groups of 5 runs. It beat the authors' average in 3 of those 4 groups. Against the retrained baseline, the merged model gains 1.89 wF1 and is better in all 20 runs. A single run sits within 0.2 wF1 of the authors, 8 of the 20 runs beat them, and the best run reaches 72.95.

\subsection{MELD: L2 + L3 (Keshav Mishra)}

This system uses CACG with no gate loss, the authors' focal loss, LLM fusion at $w = 0.2$, and temperature scaling. It was retrained on 3 seeds in \path{Merge-L2-L3/ext}, which is a copy of the rewritten code with L2 added, so unlike the IEMOCAP merge it does not run on the authors' code. Every $w$ and $T$ was fixed on dev before test was read.

\archfig{approach_a_meld}{Approach A on MELD: CACG, LLM fusion and temperature scaling}

\begin{center}
\begin{adjustbox}{max width=\linewidth}
\begin{tabular}{lcccccccc}
\toprule
Row (3 seeds) & Acc & wF1 & mF1 & Fear F1 & Disgust F1 & Joy F1 & ECE + TS & Dev mF1 \\
\midrule
Baseline & 66.83 & 65.72 & 49.57 & 23.2 & 30.5 & 64.2 & 3.49 & 54.07 \\
CB0.999 & 65.80 & 65.68 & 49.50 & 22.2 & 30.0 & 63.5 & 2.56 & 54.18 \\
CACG & 67.09 & 65.96 & 49.67 & 22.8 & 29.9 & 64.4 & 3.05 & 53.45 \\
Baseline + LLM & 68.11 & 67.01 & 53.59 & 32.3 & 38.9 & 61.6 & 3.38 & 59.96 \\
CB0.999 + LLM & 68.43 & 67.66 & 54.45 & 34.9 & 39.6 & 63.0 & 3.58 & 59.38 \\
\textbf{CACG + LLM (merged)} & \textbf{68.17} & \textbf{67.11} & \textbf{53.77} & 32.9 & 39.1 & 62.1 & 3.18 & 59.51 \\
\bottomrule
\end{tabular}
\end{adjustbox}
\end{center}

\begin{center}
\begin{adjustbox}{max width=\linewidth}
\begin{tabular}{lcccc}
\toprule
Merged system against the authors' checkpoints & Authors & Merged & Change & 95\% CI \\
\midrule
Accuracy & 67.32 & 68.17 & +0.85 & [\mn0.69, +2.38] \\
wF1 & 66.46 & 67.11 & +0.65 & [\mn0.94, +2.21] \\
\textbf{mF1} & \textbf{50.87} & \textbf{53.77} & \textbf{+2.90} & \textbf{[+0.42, +5.29]} \\
\bottomrule
\end{tabular}
\end{adjustbox}
\end{center}

Only the macro-F1 interval excludes zero. Nearly all of the gain comes from the LLM: CACG given the LLM adds +0.18 mF1 [\mn0.24, +0.62], and the authors' checkpoints plus the LLM reach 54.25. CB0.999 + LLM has the best test score, but it ranked third on dev, so it was not selected.

\begin{center}
\begin{adjustbox}{max width=\linewidth}
\begin{tabular}{lccccccc}
\toprule
Emotion F1 & Neutral & Surprise & Fear & Sadness & Joy & Disgust & Anger \\
\midrule
Authors' checkpoints & 79.79 & 58.97 & 26.12 & 40.95 & 64.95 & 31.43 & 53.89 \\
Merged system & 79.62 & 61.22 & 32.88 & 46.14 & 62.05 & 39.06 & 55.42 \\
Change & \mn0.17 & +2.25 & +6.76 & +5.19 & \mn2.90 & +7.63 & +1.53 \\
\bottomrule
\end{tabular}
\end{adjustbox}
\end{center}

\section{Approach B: one rewritten codebase}

All three limitations sit in one codebase, \path{Limitation-3/ECERC/ext/}, where every component is a command-line flag. One script trains any combination on either dataset, so each ablation row differs by flags alone. All runs were retrained on a laptop with a GeForce 940MX, and each dataset still has its own final configuration, as in the ECERC paper.

\archfig{approach_b_unified}{Approach B: one network, with a separate final route for each dataset}

\begin{center}
\begin{tabular}{L{4.2cm}L{1cm}L{4.2cm}L{6.3cm}}
\toprule
Component & From & IEMOCAP & MELD \\
\midrule
Audio features & L1 & WavLM & Release features; only IEMOCAP audio was extracted \\
VAD head + soft labels & L1 & Tested, dropped & Not tested \\
CACG cause gate & L3 & Dynamic $\tau$, no gate loss & Dynamic $\tau$, utility loss \\
Zero-shot LLM fusion & L2 & No scores available & Yes, $w = 0.2$ \\
Temperature scaling & L3 & $T$ per seed & $T = 1.64$ on the fused scores \\
Seed ensemble & L3 & 5 seeds & 3 seeds \\
\bottomrule
\end{tabular}
\end{center}

\subsection{IEMOCAP ablation}

\begin{center}
\begin{adjustbox}{max width=\linewidth}
\begin{tabular}{lccccccc}
\toprule
Variant & Seeds & wF1, single & mF1, single & ECE + TS & Ensemble wF1 & Ensemble mF1 & Ensemble ECE \\
\midrule
ECERC retrained & 5 & 69.52 $\pm$ 0.97 & 68.75 & 4.55 & 70.29 & 69.45 & 2.40 \\
+ L3 CACG & 5 & 69.77 $\pm$ 1.28 & 69.06 & 3.72 & 70.68 & 70.14 & 2.79 \\
+ L1 VAD + soft labels, IS10 audio & 3 & 70.07 $\pm$ 0.24 & 69.08 & 3.76 & 70.07 & 69.09 & 3.12 \\
+ L1 WavLM & 5 & 71.60 $\pm$ 0.32 & 70.10 & 2.94 & 71.92 & 70.30 & 2.28 \\
+ L1 WavLM + VAD + soft labels & 5 & 70.93 $\pm$ 0.87 & 69.44 & 3.05 & 71.78 & 70.28 & 2.50 \\
\textbf{+ L1 WavLM + L3 CACG (final)} & 5 & \textbf{71.92 $\pm$ 0.25} & \textbf{70.65} & 2.89 & \textbf{72.77} & \textbf{71.52} & 2.50 \\
+ L1 WavLM + VAD + soft labels + L3 & 5 & 71.46 $\pm$ 0.44 & 69.96 & 2.52 & 72.66 & 71.03 & 2.81 \\
\bottomrule
\end{tabular}
\end{adjustbox}
\end{center}

WavLM gave the largest gain and cut the seed spread from 0.97 to 0.32. VAD and soft labels together cost about 0.5 to 0.7 wF1 on top of WavLM, so they were dropped. This differs from Approach A, where soft labels were dropped first and VAD was kept. Approach B never tested VAD on its own with WavLM.

\subsection{MELD ablation}

\begin{center}
\begin{adjustbox}{max width=\linewidth}
\begin{tabular}{lccccccc}
\toprule
Variant (3 seeds, $\gamma$ = 1) & wF1, single & mF1, single & ECE raw & ECE + TS & Ensemble wF1 & Ensemble mF1 & Ensemble ECE \\
\midrule
ECERC retrained & 65.90 $\pm$ 0.28 & 49.78 & 3.62 & 2.92 & 65.99 & 50.05 & 2.73 \\
+ L3 CACG, utility loss & 65.51 $\pm$ 0.08 & 49.25 & 3.97 & 3.20 & 65.83 & 49.86 & 3.90 \\
+ L2 LLM fusion & 67.01 $\pm$ 0.22 & 53.64 & 16.81 & 3.05 & 67.31 & 54.29 & 3.22 \\
\textbf{+ L2 LLM fusion + L3 CACG (final)} & \textbf{67.23 $\pm$ 0.21} & \textbf{54.10} & 14.98 & 2.99 & \textbf{67.10} & \textbf{53.82} & 2.99 \\
\bottomrule
\end{tabular}
\end{adjustbox}
\end{center}

Fusion raised raw ECE from 3.6 to 16.8, and temperature scaling brought it back to 3.1, so the fused confidence cannot be trusted without scaling. The plain ECERC + LLM ensemble is slightly above the final one, within noise.

\subsection{What the merge involved}

\begin{itemize}
  \item \textbf{Data:} WavLM features, VAD values and annotator votes were aligned to all 7,433 IEMOCAP utterances, with none missing.
  \item \textbf{Fusion:} Keshav's cached LLM scores were fused with saved outputs, because the 32B LLM cannot run on the laptop. $w$ and $T$ were set on validation only.
  \item \textbf{Correction:} the first MELD grid used $\gamma$ = 0, but the release uses $\gamma$ = 1, as Keshav spotted. MELD was retrained at $\gamma$ = 1, and all MELD figures in this section use it.
  \item \textbf{Checks:} with new flags at their defaults, early epochs match earlier logs exactly, and the rewritten model still matches the official model with a maximum difference of 0.0.
\end{itemize}

\subsection{Demo interface}

A local page replays the test conversations of both final models. Each utterance shows the predicted emotion, calibrated confidence, a low-confidence flag and the cause weights. Flagged predictions are right 47.0\% of the time on IEMOCAP and 37.7\% on MELD, against 78.7\% and 76.3\% for the rest. It cannot take typed sentences, because the final models need raw audio features, emotion-tuned RoBERTa weights we do not have, and the 32B LLM.

\section{Comparing both approaches}

The two approaches reach almost the same results by the same route, so the conclusion does not depend on which code was used. On IEMOCAP, Approach B scores 0.34 wF1 higher, and on MELD the two are level, but neither gap is proven.

\begin{center}
\begin{adjustbox}{max width=\linewidth}
\begin{tabular}{llcccc}
\toprule
IEMOCAP & Runs & Acc & wF1 & mF1 & ECE + TS \\
\midrule
Authors' checkpoints & 3 & 71.60 & 71.78 & 70.82 & 2.81 \\
Approach A, single run & 20 & 71.36 & 71.59 & 70.19 & 3.41 \\
Approach B, single run & 5 & 71.77 & 71.92 & 70.65 & 2.89 \\
Approach A, ensemble & 5-run, 4 groups & 72.18 & 72.43 & 71.14 & 2.12 \\
Approach B, ensemble & 5-run, 1 group & 72.58 & 72.77 & 71.52 & 2.50 \\
\bottomrule
\end{tabular}
\end{adjustbox}
\end{center}

\begin{center}
\begin{adjustbox}{max width=\linewidth}
\begin{tabular}{llcccc}
\toprule
MELD & Runs & Acc & wF1 & mF1 & ECE + TS \\
\midrule
Authors' checkpoints & 3 & 67.32 & 66.46 & 50.87 & 3.02 \\
Authors' checkpoints + LLM & 3 & 68.24 & 67.26 & 54.25 & \\
Approach A, mean over seeds & 3 & 68.17 & 67.11 & 53.77 & 3.18 \\
Approach B, mean over seeds & 3 & 68.28 & 67.23 & 54.10 & 2.99 \\
Approach B, ensemble & 3 & 68.16 & 67.10 & 53.82 & 2.99 \\
\bottomrule
\end{tabular}
\end{adjustbox}
\end{center}

\subsection{Where the approaches differ}

{\small
\begin{longtable}{L{2.4cm}L{3.6cm}L{3.6cm}L{6.2cm}}
\toprule
Choice & Approach A & Approach B & What it means \\
\midrule
\endhead
\bottomrule
\endlastfoot
IEMOCAP code & Authors' code & Rewritten code & Same model, so the code is not the main difference \\
MELD code & Copy of the rewritten code & Rewritten code & Both MELD systems use the same core model code \\
IEMOCAP loss & WavLM + VAD & WavLM only & VAD adds +0.15 wF1 in Approach A, within noise \\
MELD gate loss & None & Utility loss & Scores are level, but weights are only readable without the loss \\
IEMOCAP seeds & 20 & 5 & Approach A's evidence is much stronger \\
Model selection & Validation, with VAD kept after comparing scores & Test ablation for the IEMOCAP choice & Approach B's IEMOCAP result is slightly optimistic \\
Coverage & One dataset per system & Both datasets, with a demo & Approach B is simpler to run and show \\
\end{longtable}
}

\subsection{What both approaches agree on}

\begin{itemize}
  \item \textbf{New evidence beat new losses.} WavLM on IEMOCAP and the zero-shot LLM on MELD gave the gains, after loss and sampling changes had failed in both limitations.
  \item \textbf{CACG adds no proven accuracy.} Its value is readable cause weights and a small calibration gain.
  \item \textbf{Temperature scaling is needed.} It keeps confidence honest, above all after LLM fusion.
  \item \textbf{Calibration ties with the authors once scaled.} After the same scaling, the authors' checkpoints reach 2.81 and 3.02, about level with our final systems. Most of the drop from raw ECE comes from scaling, not from our models.
\end{itemize}

\subsection{Final approach}

The final method is one recipe with a different evidence source per dataset. On IEMOCAP, it is ECERC with WavLM audio, CACG, temperature scaling and a 5-seed ensemble. On MELD, it is ECERC with CACG, zero-shot LLM fusion and temperature scaling, and CACG should run without the utility loss, so its weights stay readable.

For headline numbers, we recommend quoting Approach A's 72.43 wF1 on IEMOCAP as the best-evidenced result, and Approach B's 72.77 as the matching result in the unified code.

\section{Individual contributions}

Each member owned one limitation, and all three reproduced ECERC on their datasets.

{\small
\begin{longtable}{L{3.2cm}L{8.1cm}L{4.8cm}}
\toprule
Member & Main work & Key result \\
\midrule
\endhead
\bottomrule
\endlastfoot
\textbf{Divyansh Dubey} (12340690) & L1 on IEMOCAP: six methods, WavLM extraction for 7,433 utterances, VAD and vote labels, 20-seed paired study, L1 + L3 merge on the authors' code, IEMOCAP final report & wF1 +1.43; 20\% fewer clear-cut Angry--Frustrated errors; ensemble 72.43 wF1 \\
\textbf{Keshav Mishra} (12341140) & L2 on MELD: 12 methods, pre-registered zero-shot LLM fusion, L2 + L3 merge, MELD final report and architecture figures. Spotted the $\gamma$ = 1 loss & mF1 +3.38 from fusion; merged 53.77 mF1, +2.90 over the authors \\
\textbf{Lakshay Gupta} (12341300) & L3 on both datasets, plus the rewritten code and the all-in-one merge (details below) & CACG, ECE 14.36 to 3.72 with scaling, ensemble 72.77 wF1 and 53.82 mF1 \\
\end{longtable}
}

\textbf{Lakshay Gupta's work in detail:}

\begin{itemize}
  \item \textbf{Unified code:} rewrote ECERC as one package for both datasets, with a maximum output difference of 0.0 from the official model. Teammates' MELD merge was built on it.
  \item \textbf{Code audit:} found the missing cause softmax, the empty-history attention problem and the MELD modality routing.
  \item \textbf{CACG:} designed the masked cause distribution, 12-d descriptor, per-utterance temperature and availability mask.
  \item \textbf{Utility gate loss:} designed a label-grounded replacement for the proposal's entropy loss.
  \item \textbf{Interpretability:} exposed per-utterance cause weights, and showed which causes the model relies on for each emotion.
  \item \textbf{Calibration:} added ECE, NLL and reliability analysis, temperature scaling and cause-aware scaling.
  \item \textbf{Seed ensembles and selective prediction:} turned confidence into a low-confidence flag.
  \item \textbf{Ablation:} ran the 42-run L3 grid and the merged ablation on both datasets.
  \item \textbf{Approach B:} merged all three limitations behind flags, retrained on one laptop, and built the demo interface.
  \item \textbf{Documentation:} wrote \path{FINAL_MODEL.md}, a merging guide for teammates, the SVG architecture diagrams and this combined report.
\end{itemize}

\section{Limits of these results}

The gains are real but modest in certainty, and several caveats apply to how they should be read.

\begin{itemize}
  \item \textbf{Repeated test use.} The MELD test set has been used four or more times. After correcting for four looks, the interval on the +2.90 mF1 gain is [\mn0.30, +5.96], which includes zero.
  \item \textbf{Selection.} Approach B chose its IEMOCAP model from test scores, because validation preferred the original IS10 audio. On MELD, dev ranks ECERC + LLM above the merged system, which was named because it combines both limitations.
  \item \textbf{Seeds.} Approach B used 5 and 3 seeds, so small differences cannot be confirmed. Only Approach A's IEMOCAP study has 20 paired seeds.
  \item \textbf{Small classes.} Test has only 50 Fear and 68 Disgust utterances, so per-emotion changes are noisy.
  \item \textbf{Baseline variation.} MELD retraining gave 65.72 to 66.00 wF1 across three runs of the same recipe, and the cause is not established.
  \item \textbf{Inference cost.} The MELD system needs a 32B LLM, used here through cached scores. WavLM needs raw audio.
  \item \textbf{LLM exposure.} The LLM's pretraining data may include Friends transcripts or MELD's public files. Alone it is weaker (63.14 wF1), which argues against memorised labels.
  \item \textbf{No cross-dataset tests.} WavLM was not tested on MELD, and LLM fusion was not tested on IEMOCAP.
  \item \textbf{Code availability.} Approach A's IEMOCAP merge code is in \path{Merge-L1-L3/}, but its run outputs are not in the repository, so its figures come from the written report.
\end{itemize}

\section{Sources}

\begin{itemize}
  \item ECERC: T. Zhang and Z. Tan, \emph{ECERC: Evidence-Cause Attention Network for Multi-Modal Emotion Recognition in Conversation}, ACL 2025.
  \item Limitation 1 and IEMOCAP merge: \texttt{ECERC Revisited on IEMOCAP Final Report (Limitations 1 and 3).pdf}, \path{Merge-L1-L3/README.md}.
  \item Limitation 2 and MELD merge: \path{Merge-L2-L3/REPORT.md}, \path{Merge-L2-L3/README.md}, \path{Limitation-2/README.md}.
  \item Limitation 3 and Approach B: \path{FINAL_MODEL.md}, \path{Limitation-3/ECERC/runs/RESULTS.md}, \path{Limitation-3/ECERC/runs/RESULTS_merged.md}.
  \item Architecture diagrams: the SVG files in \path{architecture/}, generated by \path{architecture/make_svgs.py}.
\end{itemize}

\end{document}
